\subsection{Analysis of Muskingum Parameters}
\label{subsec:analysis-muskingum-parameters}

The Muskingum equation is given in Equation~\ref{eq:muskingum} \citep{usace_hms} where $Q$ and $I$ are channel flows in cubic meters per second
at the indicated time step and river reach. The coefficients are given in Equations~\ref{eq:c1},~\ref{eq:c2}, and~\ref{eq:c3}. The Muskingum
equation's three coefficients, $c_1$, $c_2$, and $c_3$, are functions of the Muskingum parameters $k$ and $x$, and the time step $\Delta t$. The
Muskingum equation is stable, meaning that errors are not generated or amplified by the equation, for any $x$ between 0.0 and 0.5
\citep{mccarthy1938}. Valid values of $k$ and $\Delta t$ are positive numbers.

\begin{equation}
    Q_{t+1} = c_1\, I_{t+1} + c_2\, I_t + c_3\, Q_t
    \label{eq:muskingum}
\end{equation}

\begin{equation}
    c_1 = \frac{\Delta t / k - 2x}{\Delta t / k + 2 (1-x)}
    \label{eq:c1}
\end{equation}

\begin{equation}
    c_2 = \frac{\Delta t / k + 2x}{\Delta t / k + 2 (1-x)}
    \label{eq:c2}
\end{equation}

\begin{equation}
    c_3 = \frac{2 (1-x) - \Delta t / k}{\Delta t / k + 2 (1-x)}
    \label{eq:c3}
\end{equation}

Muskingum only produces physically valid results when the coefficients are positive \citep{hjelmfelt1985, chow1988, gill1992}. These conditions
are summarized in Equation~\ref{eq:validity-conditions}. The three coefficients share a common denominator which is positive for all valid values
of $k$, $x$, and $\Delta t$. The numerator of $c_2$ is always positive for valid parameter values reducing the validity condition to values which
produce positive numerators of $c_1$ and $c_3$. The validity condition is summarized in Equation~\ref{eq:validity-conditions}. When $x$is at the
maximum of 0.5, meaning no attenuation, $\Delta t$ and $k$ must be set equal. When $x$ is at the minimum of 0.0, meaning maximum attenuation,
$\Delta t$ must be less than $2k$.

\begin{gather}
    c1: \Delta t > 2 k x \notag \\
    c3: \Delta t < 2 k (1-x) \notag \\
    2 k x < \Delta t < 2 k (1-x)
    \label{eq:validity-conditions}
\end{gather}

The two primary cases where negative coefficients are created are first $k > \Delta t$ causing $c_1$ to be negative, and second when $\Delta t > k$
causing $c_3$ to be negative. When $c_1$ is negative, the inflow at the next time step $I_{t+1}$ is subtracted from the outflow $Q_{t+1}$. This is
incorrect because larger inflows should produce larger outflows given the assumptions of wedge and prism storage in Muskingum routing. This phenomenon
is documented in % todo: add citations.
When $c_3$ is negative, the outflow at the previous time step $Q_t$ is subtracted from the outflow at the next time step $Q_{t+1}$.
A negative $c_3$ can produce negative discharge on $Q_{t+1}$ if $Q_t$ is larger than the inflow terms $I_{t+1}$ and $I_t$. Subsequent time steps then
weight the negative $Q_{t}$ with negative $c_3$ switching back to a positive $Q_{t+1}$. The oscillation dampens over time because the solution is
stable. Oscillating flow has been documented in % todo: add citations.

\subsection{Mitigation Strategies for Invalid Parameterization}
\label{subsec:mitigation-strategies}

%    todo: now talk about the mitigation strategies
%    todo: Show how the Courant number is related to the validity conditions and how it can be used to determine if simulation is valid or not.
%    todo: paragraph about stability and validity relating to the Courant number, 3 coefficients dt, k, x, and dx

\subsection{Analysis of Global Hydrography Products}
\label{subsec:analysis-hydrography-products}

We surveyed a total of six global hydrography products assessing their fitness for routing calculations. The datasets are HydroRIVERS (HydroSHEDS v1)
\citep{lehner2013},HydroSHEDS~v2 \citep{lehner2022, lehner2026}, MERIT-Basins \citep{lin2019}, NHDPlus~V2 \citep{mckay2012}, TDX-Hydro
\citep{carlson2024}, and GRIT v1 \citep{wortmann_grit, wortmann2025}. Four of these products have global coverage. The three most recently produced
datasets, HydroSHEDS v2, TDX-Hydro, and GRIT, are derived from the 12-meter TanDEM-X elevations or the FABDEM 30-meter elevations derived from
TanDEM-X. Two, HydroRIVERS and MERIT-Basins, ultimately derive from the 90-meter SRTM elevations. NHDPlus~v2 comes from USGS 3DEP elevation data which
is the only non-global elevation product represented.

Streams were analyzed primarily for the distribution of their reach lengths and their $k$ parameters, calculated as $k = Length / velocity$ assuming
constant 0.5 m/s velocity. Using a fixed velocity for all streams keeps the distribution of $k$ identical to the distribution of lengths and functions
as a unit converter. The baseline speed of 0.5 m/s is used a reasonable starting point for exploration and should be interpreted or conflated with
actual velocity distributions or calibrated velocities.

\subsection{Analysis of Existing Routing Code Implementations}
\label{subsec:analysis-of-routing-code}

This article analyzes Muskingum implementations of WRF-Hydro as extracted in t-route \citep{troute_code} and RAPID \citep{david2011}.
and the other analysis of this paper provide the analytical basis for creating the river-route implementation intended for future
RFS models.
